\documentclass[10pt,a4paper]{article}
\PassOptionsToPackage{table}{xcolor}
\usepackage{ocad}
\graphicspath{{charts/}}
\hypersetup{pdftitle={OCAD - Banco de Chile},pdfauthor={Arthur Faugeron}}
\ocadSetReportID{OCAD report: Banco de Chile (BCH), 1 November 2026}

\begin{document}

\ocadcoversetup{company={Banco de Chile}, ticker={BCH}, exchange={NYSE (ADS); Santiago Stock Exchange (CHILE)},
  series={Bank}, date={1 November 2026}, analyst={Arthur Faugeron},
  price={USD 36.55 per ADS}, pricenote={Closing price, Friday 30 October 2026, latest completed trading day},
  chart={cover_price.png}, chartcaption={BCH daily closing price, 3 August to 30 October 2026, USD per ADS. Source: Alpha Vantage daily series (NYSE).}}
\makeocadcover

\section{Executive Summary}
\ocadlede{Banco de Chile is the second-largest private bank in Chile by loans, controlled by the Quiñenco and Citigroup joint holding LQ Inversiones Financieras. It earns a 21 percent return on average equity on a 37 percent efficiency ratio, but the first-half 2026 margin is flattered by inflation and management itself guides to a lower full-year margin. The ADS closed at USD 36.55, its lowest close in the last 100 sessions, after falling 4.5 percent on the final trading day of October.}

\subsection{Business}
Banco de Chile is a universal commercial bank serving retail customers, small and medium enterprises and large corporates, plus treasury, insurance brokerage and investment products, with more than 2.9 million active clients and about 11,000 employees. Net interest income was 57.7 percent of 2025 operating revenue before loan losses and non-interest income 42.3 percent, all in Chile apart from a representative office in Beijing. It reports under IFRS in Chilean pesos (CLP), and each ADS represents 200 ordinary shares.

\subsection{Financial Condition}
FY2025 operating revenue before loan losses was CLP 3,026 bn (down 0.8 percent), net income CLP 1,192 bn (down 1.3 percent), return on average equity 21.2 percent as reported, and the equity to assets ratio 10.7 percent. First half 2026 net income was CLP 659 bn (up 4.0 percent), with second quarter net income of CLP 391 bn (up 28.1 percent) helped by a net interest margin of 5.76 percent. Bank cash flow from operations is dominated by changes in trading securities and loans, so free cash flow is not a reliable measure for a bank; dividends are the relevant cash measure.

\subsection{Market Position}
The main competitors are Banco Santander-Chile, BCI, BancoEstado, Itaú Chile and Scotiabank Chile. Its major advantage is a low-cost funding base, with non-interest bearing deposits at 41 percent of deposits, and a 37 percent efficiency ratio. Its major weakness is dependence on one economy, with 2026 growth guided near 1.3 percent, and on inflation-linked income.

\subsection{Valuation}
At USD 36.55 the market capitalization is USD 18.5 bn. Trailing P/E is 15.2 times and forward P/E about 13.7 times, price to book 3.1 times, and the dividend yield 4.5 percent (all rescaled to the 30 October close). Enterprise value, EV/EBITDA and P/FCF are not meaningful for a deposit-funded bank.

\subsection{Investment Framework}
\begin{ocadframework}
\fwrow{What is attractive}{Return on equity above 20 percent for four consecutive years, a cost to income ratio near 37 percent, no share dilution, and a dividend yield of 4.5 percent.}
\fwrow{What is concerning}{Net interest margin is guided to 4.6 percent for 2026 against 4.93 percent in the first half, credit loss expense rose 49.7 percent in the first half, and the stock trades at 3.1 times book.}
\fwrow{What must go right}{Loan growth near 6 percent, credit cost within 1.2 to 1.3 percent of average loans, and a return on average capital of 21 to 22 percent, as guided.}
\fwrow{What could break the thesis}{Faster margin normalization as inflation and policy rates fall, a rise in consumer and SME arrears, or a sustained return on equity below 17 percent.}
\fwrow{Priced in}{At a 10 percent cost of equity and 3.5 percent growth, a price to book of 3.1 times requires a perpetual return on equity near 23.5 percent, above the 19 to 21 percent currently earned.}
\end{ocadframework}
\needspace{16\baselineskip}
\subsection{Investment Profile}
\begin{ocadtable}{Assessment by dimension}{>{\raggedright\arraybackslash}p{40mm} L}
\ocadtablehead{\ocadH{Dimension}&\ocadH{Assessment}}
Business quality & High. Scale, a 2.9 million client base and a low-cost deposit franchise.\ocadhr
Financial quality & High. Return on equity of 20.9 percent and 37 percent efficiency, with earnings falling since 2022.\ocadhr
Competitive position & Strong. Leading private-bank franchise in a concentrated market.\ocadhr
Growth profile & Low. Revenue CAGR of 8.0 percent since 2021, but flat since 2022; 2026 loan growth guided near 6 percent.\ocadhr
Balance-sheet risk & Low to moderate. Equity to assets 10.7 percent; regulatory capital ratios not retrieved.\ocadhr
Valuation & Full. 3.1 times book and 15.2 times earnings, a premium to local peers.\ocadhr
Market and technical profile & Weak near term. Closing at a 100-session low, below the 20 and 50-day averages, RSI near 30.\ocadhr
\end{ocadtable}
\ocadtablenote{Source: analyst assessment based on the data in this report.}

\section{A. About the Business}
\subsection{A1. Business Environment}
The business is commercial banking in Chile. Banks gather deposits and wholesale funding, lend to households, firms and the state sector, and earn the spread plus fees, with credit losses and operating costs as the main deductions. According to the second quarter 2026 filing, system loans were CLP 253,069,875 million at 31 May 2026, up 3.8 percent in nominal terms. With 2026 inflation near 4.0 percent in management's assumptions, real loan growth is therefore close to zero. Value is created by low-cost funding and cheap distribution, and captured by the banks with scale: Banco de Chile's 2025 cost to income ratio was 37.4 percent as reported.

The macro backdrop is moderate. Chile's GDP grew 0.5 percent in 2023 and 2.6 percent in 2024, with 2.4 percent expected for 2025, while consumer price inflation was 3.9, 4.5 and 3.4 percent and the Central Bank policy rate ended each year at 8.25, 5.00 and 4.50 percent. Management assumes 2026 GDP growth near 1.3 percent, December inflation near 4.0 percent and a policy rate near 4.5 percent. The industry is mature, cyclical and highly regulated: growth tracks nominal GDP, and profitability swings with inflation and interest rates because many loans and securities are indexed to the unit of account (UF).
\ocadchart{mcap.png}{Banco de Chile market capitalization at year end, 2012 to 2025, and 30 October 2026}{Source: StockAnalysis.com market capitalization history; Alpha Vantage close of 30 October 2026. Units: USD bn.}
\begin{ocadcallout}{Industry stage: mature, cyclical and regulated}
Loan growth near nominal GDP, concentrated competition among six large banks, and earnings that depend on inflation and policy rates indicate a mature industry. Digital channels and open banking are structural changes, but the retrieved sources do not quantify their effect on margins.
\end{ocadcallout}

\subsection{A2. What Does the Company Do?}
Banco de Chile takes deposits and issues bonds in pesos, UF and foreign currency, lends the proceeds to individuals, SMEs and corporates, runs a treasury that trades securities and manages liquidity, and sells insurance and investment products through its subsidiaries. A new employee should think of three engines: the spread on a CLP 39 trillion loan book, the zero-cost funding of CLP 11.8 trillion of current account deposits, and fee income, which was CLP 340 bn in the first half of 2026 (up 8.8 percent). The loan book was CLP 39.2 trillion at the end of 2025.
\subsubsection{How one peso of operating revenue is used}
Using 2025 figures, each CLP 100 of operating revenue before loan losses (CLP 3,026 bn) is distributed as follows: CLP 12.6 to provisions for loan losses, CLP 37.3 to operating expenses, CLP 5.4 to other items before tax, CLP 10.6 to income tax, and CLP 39.4 to net income. Of net income, roughly two thirds is distributed as dividends.
\begin{ocadflow}{3}
\flowbox{1. Customer}{Individuals, SMEs, corporates; 2.9 million active clients}
\flowbox{2. Product and service}{Loans, deposits, cards, treasury, insurance, investments}
\flowbox{3. Revenue}{CLP 1,746 bn net interest income and CLP 1,280 bn non-interest income}
\end{ocadflow}
\begin{ocadflow}{4}
\flowbox{4. Costs}{CLP 382 bn provisions; CLP 1,130 bn operating expenses}
\flowbox{5. Operating profit}{CLP 1,513 bn pre-tax; 50.0 percent of revenue}
\flowbox{6. Cash flow}{Net income CLP 1,192 bn; bank cash flow is distorted by trading assets}
\flowbox{7. Capital allocation}{Dividends of about two thirds of earnings; no buybacks; retained earnings build equity}
\end{ocadflow}
\ocadtablenote{Source: Banco de Chile FY2025 income statement as compiled by StockAnalysis.com, CLP bn. Distribution calculated from reported lines.}

\subsection{A3. Business and Relationship Network}
\begin{ocadtable}{Ecosystem map}{>{\raggedright\arraybackslash}p{30mm} L L}
\ocadtablehead{\ocadH{Node}&\ocadH{Relationship}&\ocadH{Why it matters economically}}
LQ Inversiones Financieras (LQIF) & Holds 51.15 percent of the shares as of 30 June 2026; owned equally by Quiñenco and Citigroup since 2010 & Control is stable and the controlling holders have held the position since 1999, but control means a thin free float.\ocadhr
Quiñenco & Chilean holding company; carries Banco de Chile at CLP 1,894.6 bn book value & Largest asset in the group; dividends from the bank fund the holding.\ocadhr
Citigroup & Partner in LQIF since 2008; international connectivity alliance & Supports cross-border corporate business and trade finance.\ocadhr
Depositors and bondholders & CLP 28.5 trillion of deposits and CLP 16.2 trillion of debt at end 2025 & Funding cost drives the spread; 41 percent of deposits pay no interest.\ocadhr
Borrowers & Households, SMEs and corporates in Chile & Credit cost and loan growth depend on the domestic cycle.\ocadhr
Regulators & Central Bank of Chile and the financial markets commission (CMF) & Set policy rates, capital and liquidity rules.\ocadhr
Competitors & Santander-Chile, BCI, BancoEstado, Itaú, Scotiabank & Share loans, deposits and talent in a concentrated market.\ocadhr
\end{ocadtable}
\ocadtablenote{Source: Quiñenco, Banco de Chile company page; Banco de Chile Form 6-K, 2Q26 Financial Management Review; StockAnalysis.com balance sheet.}

\subsection{A4. Competition}
\begin{ocadtable}{Competitive comparison}{>{\raggedright\arraybackslash}p{34mm} R R R R R}
\ocadtablehead{\ocadH{Bank}&\ocadH{Mkt cap, USD bn}&\ocadH{P/E, x}&\ocadH{P/B, x}&\ocadH{ROE, \%}&\ocadH{Div. yield, \%}}
Banco de Chile (BCH) & 18.5 & 15.2 & 3.08 & 21.1 & 4.5\ocadhr
Banco Santander-Chile (BSAC) & 14.4 & 11.8 & 2.58 & 23.6 & 3.6\ocadhr
Bancolombia (CIB) & 24.6 & 16.2 & 2.14 & 19.2 & 4.9\ocadhr
Credicorp (BAP) & 29.7 & 13.8 & 2.66 & 20.4 & 4.0\ocadhr
\end{ocadtable}
\ocadtablenote{Source: StockAnalysis.com statistics pages. BCH rescaled to the 30 October 2026 close and BSAC from 24 September to 30 October 2026 by the change in price; CIB (26 October 2026) and BAP as published, unscaled. P/E is trailing; ROE as published.}
Banco Santander-Chile is the only listed direct competitor with a similar corporate and retail mix and is the main valuation comparison. BCI, BancoEstado (the main state-owned bank), Itaú Chile and Scotiabank Chile compete for the same customers, but per-bank figures were not retrieved and are not quoted. Bancolombia and Credicorp are regional listed banks used only to place the valuation in a Latin American context; they do not compete for Chilean customers. Fintech lenders and card issuers are emerging substitutes for payments and consumer credit. Santander-Chile reported FY2025 net income of CLP 1,022 bn against Banco de Chile's CLP 1,192 bn, and retail customers were 76.3 percent of its loan book.
\ocadchart{peers.png}{Valuation and returns versus listed banks}{Source: StockAnalysis.com, late October 2026 (see table note for dates). Units: multiples in times, ROE in percent.}

\subsection{A5. Competitive Differences}
Advantages are mostly structural: a funding franchise in which CLP 11.8 trillion of deposits carry no interest, a 2.9 million client base, and an operating cost base that produced a 37.3 percent calculated efficiency ratio in 2025. The 51.15 percent controlling holder provides capital stability. Disadvantages are concentration in one economy, a margin sensitive to inflation and policy rates, and a loan book where growth is limited by the system's 3.8 percent nominal expansion. Brand and scale are durable; the funding advantage is partly cyclical, since the share of non-interest bearing deposits fell from 56.7 percent in 2021 to 41.4 percent in 2025 as rates rose.
\begin{ocadcallout}{Moat assessment: Moderate to Strong}
Return on equity of 20.9 to 29.7 percent over four years exceeds any plausible cost of equity, which supports a durable franchise, but part of it reflects the 2022 inflation spike and the loss of zero-cost funding shows the advantage is not fixed.
\end{ocadcallout}

\subsection{A6. Industry Trend}
System loan growth of 3.8 percent nominal is a mature, cyclical trend. Structural drivers are digitalization and regulatory capital requirements; temporary conditions are the high inflation of 2021 to 2023 and the cycle of policy rates, which fell from 8.25 percent to 4.50 percent between 2023 and 2025. The trend could reverse upward with faster GDP growth or downward if arrears rise: the industry past-due ratio is forecast at 2.3 to 2.4 percent for 2026 in the company's filing.

\subsection{A7. Management}
\begin{ocadtable}{Executive team}{>{\raggedright\arraybackslash}p{36mm} >{\raggedright\arraybackslash}p{44mm} L}
\ocadtablehead{\ocadH{Name}&\ocadH{Position}&\ocadH{Notes}}
Eduardo Ebensperger Orrego & Chief Executive Officer (General Manager) & Leads the bank; further biography not retrieved.\ocadhr
Rolando Arias Sanchez & Chief Financial Officer & Responsible for finance; biography not retrieved.\ocadhr
Hernan Buchi Buc & Director & Board member listed by the data provider.\ocadhr
Pablo Camilo Mejia Ricci & Head of Investor Relations & Main contact for the 6-K and 20-F disclosures.\ocadhr
Alfredo Villegas Montes & General Counsel and Secretary of the Board & Corporate governance and legal.\ocadhr
Rodrigo Aravena & Chief Economist and Senior Vice President of Institutional Relations & Source of the macro assumptions.\ocadhr
\end{ocadtable}
\ocadtablenote{Source: StockAnalysis.com company profile, October 2026. Previous careers and compensation were not retrieved.}
Management behavior is shareholder-aligned on distributions and conservative on structure. The share count has been 101,017 million for every year from 2021 to 2025, so there has been no dilution and no buyback. Dividend yield of about 4.5 percent and a published payout near 67 percent return most of the earnings, while retained earnings rose from CLP 1,124 bn to CLP 2,677 bn between 2021 and 2025. Guidance accuracy is mixed: the 2026 loan growth target was lowered from about 7 to about 6 percent, whereas first half return on average capital of 22.4 percent is already at the top of the 21 to 22 percent range. Goodwill was nil in 2025, so there is no acquisition risk on the balance sheet. The controlling holders are the insiders: Quiñenco and Citigroup hold 51.15 percent through LQIF.

\section{B. Evaluation}
\subsection{B1. Financial and Operating Metrics}
\begin{ocadtable}{Profitability, efficiency and credit quality}{>{\raggedright\arraybackslash}p{50mm} R R R R R}
\ocadtablehead{\ocadH{Metric}&\ocadH{FY2021}&\ocadH{FY2022}&\ocadH{FY2023}&\ocadH{FY2024}&\ocadH{FY2025}}
Operating revenue before loan losses, CLP bn & 2,224.9 & 3,142.6 & 2,994.9 & 3,050.3 & 3,026.0\ocadhr
Revenue growth, \% & N/A & 41.3 & -4.7 & 1.8 & -0.8\ocadhr
Pre-tax margin on revenue, \% & 59.9 & 55.2 & 50.7 & 50.0 & 50.0\ocadhr
Net margin on revenue, \% & 47.5 & 46.0 & 41.5 & 39.6 & 39.4\ocadhr
ROE on average equity, \% & N/A & 29.7 & 23.3 & 22.2 & 20.9\ocadhr
ROA on average assets, \% & N/A & 2.71 & 2.24 & 2.24 & 2.25\ocadhr
Efficiency ratio, \% & 39.2 & 31.7 & 37.2 & 37.1 & 37.3\ocadhr
Provisions to gross loans, \% & 0.06 & 1.10 & 0.94 & 0.98 & 0.96\ocadhr
Effective tax rate, \% & 20.7 & 16.7 & 18.1 & 20.9 & 21.2\ocadhr
Equity to assets, \% & 8.3 & 9.9 & 9.4 & 10.8 & 10.7\ocadhr
Diluted EPS, CLP per share & 10.46 & 14.31 & 12.31 & 11.95 & 11.80\ocadhr
\end{ocadtable}
\ocadtablenote{Source: Banco de Chile financial statements as compiled by StockAnalysis.com (reported); ratios calculated by the analyst. Efficiency ratio is non-interest expense over operating revenue before loan losses. Gross margin, ROIC, EBITDA and inventory metrics are not meaningful for a bank and are omitted.}
Revenue grew at a 8.0 percent CAGR from 2021 to 2025 but has been flat since 2022, while net income CAGR over 2021 to 2025 was 3.1 percent and over 2022 to 2025 was negative 6.2 percent. Return on equity has compressed from 29.7 to 20.9 percent as the 2022 inflation windfall unwound and equity grew; net interest income fell 22.9 percent in 2023 and recovered 15.2 percent in 2024. A cost to income ratio stable at 37 percent, provisions near 1.0 percent of gross loans and a 21 percent tax rate show a steady operating model whose swing factor is the margin.
\ocadchart{revni.png}{Operating revenue before loan losses and net income, FY2021 to FY2025}{Source: Banco de Chile financial statements as compiled by StockAnalysis.com. Units: CLP bn.}
\ocadchart{ratios.png}{Return on equity, return on assets and efficiency ratio, FY2021 to FY2025}{Source: calculated from StockAnalysis.com financial statements. Units: percent.}

\subsection{B2. Risks}
\begin{ocadtable}{Risk register ranked by probability times impact}{>{\raggedright\arraybackslash}p{9mm} >{\raggedright\arraybackslash}p{34mm} >{\raggedright\arraybackslash}p{17mm} >{\raggedright\arraybackslash}p{16mm} L}
\ocadtablehead{\ocadH{Rank}&\ocadH{Risk}&\ocadH{Probability}&\ocadH{Impact}&\ocadH{Description}}
1 & Margin normalization & High & High & Guidance of 4.6 percent for 2026 compares with 4.93 percent in the first half and 5.76 percent in the second quarter; lower inflation cuts UF indexed income.\ocadhr
2 & Credit cost increase & Medium & High & Credit loss expense rose 49.7 percent in the first half; the second quarter included CLP 50 bn of additional provisions.\ocadhr
3 & Valuation compression & Medium & High & At 3.1 times book the stock needs a return on equity above 23 percent to justify the price at a 10 percent cost of equity.\ocadhr
4 & Domestic macro and policy & Medium & Medium & GDP growth near 1.3 percent in 2026; policy rate near 4.5 percent.\ocadhr
5 & Funding mix & Medium & Medium & Non-interest bearing deposits fell from 56.7 to 41.4 percent of deposits between 2021 and 2025.\ocadhr
6 & Regulation and capital & Low & Medium & Capital and liquidity ratios were not retrieved; capital rules could limit payouts.\ocadhr
7 & Technology and competition & Medium & Low & Digital entrants pressure payments and consumer credit fees.\ocadhr
8 & Concentrated control & Low & Medium & 51.15 percent held by LQIF; sale or governance change would affect the free float.\ocadhr
\end{ocadtable}
\ocadtablenote{Source: analyst assessment using Banco de Chile Form 6-K, 2Q26, and StockAnalysis.com balance sheet.}
The three risks most capable of permanently impairing shareholder value are margin normalization, a credit cycle that raises loan losses well above 1.3 percent of loans, and valuation compression from a price that already embeds a return on equity above current earnings.

\subsection{B3. Valuation}
Using the 30 October close of USD 36.55: market capitalization is USD 18.46 bn; trailing P/E is 15.2 times (published 15.9 times at 38.33); forward P/E about 13.7 times (published 14.35 times at 38.33); price to book 3.08 times; price to tangible book 3.2 times; and the dividend yield is 4.53 percent on the USD 1.65471 per ADS paid in April 2026. P/S, EV/EBITDA, EV/FCF and PEG are not meaningful for a deposit-funded bank. The tangible book value per ADS is implied by the published price to book and equals about USD 11.85.

Against peers, BCH carries the highest price to book of the four banks in the comparison at 3.08 times against 2.58 times for Banco Santander-Chile, though its ROE of 21.1 percent is below Santander-Chile's 23.6 percent. Historical multiple ranges were not retrieved; the year-end market capitalization has ranged from USD 7.9 bn in 2021 to USD 19.2 bn in 2025.

\textbf{What is priced in.} The price to book ratio equals (ROE minus g) divided by (cost of equity minus g). At a 10 percent cost of equity and 3.5 percent perpetual growth, 3.08 times book requires a perpetual ROE of 23.5 percent. Management guides to a 21 to 22 percent return on average capital in 2026 and the calculated 2025 ROE is 20.9 percent, so expectations appear optimistic: they assume return on equity holds at the cycle peak. The cost of equity is an analyst assumption because the Chilean long-term yield was not retrieved.

\subsection{B4. Capital Market and Stock Analysis}
\begin{ocadtable}{Price returns, BCH ADS}{>{\raggedright\arraybackslash}p{58mm} R L}
\ocadtablehead{\ocadH{Period}&\ocadH{Return}&\ocadH{Basis}}
1 day (30 October 2026) & -4.5\% & Calculated, daily closes\ocadhr
1 week (5 sessions) & -4.6\% & Calculated\ocadhr
1 month (21 sessions) & -6.2\% & Calculated\ocadhr
3 months (63 sessions) & -12.0\% & Calculated\ocadhr
Since 31 December 2025 & -3.8\% & Derived: year-end market cap divided by ADS count\ocadhr
1 year to 23 October 2026 & +15.1\% & Published by StockAnalysis.com\ocadhr
Since year end 2023 & +58.3\% & Derived, price only\ocadhr
Since year end 2021 & +132.8\% & Derived, price only\ocadhr
\end{ocadtable}
\ocadtablenote{Source: Alpha Vantage daily series; StockAnalysis.com. Returns exclude dividends. Year-end prices are market capitalization divided by 505.085 million ADSs. 6 month and 3 year returns were not retrieved.}
\ocadchart{tech.png}{BCH price with 20 and 50-day averages and volume, 10 June to 30 October 2026}{Source: Alpha Vantage daily series. Units: USD per ADS; volume in thousands of ADSs.}
\begin{ocadtable}{Volatility, risk and technical indicators}{>{\raggedright\arraybackslash}p{58mm} R L}
\ocadtablehead{\ocadH{Measure}&\ocadH{Value}&\ocadH{Basis}}
Annualized volatility (100 sessions) & 25.5\% & Calculated from daily log returns\ocadhr
Downside volatility (annualized) & 18.9\% & Calculated\ocadhr
Maximum drawdown in the window & -15.1\% & High of USD 43.04 on 24 August to USD 36.55\ocadhr
Beta, five years & 0.06 & Published by StockAnalysis.com\ocadhr
Beta and correlation versus BSAC & 0.97 and 0.88 & Calculated, 99 daily returns\ocadhr
RSI (14 days) & 29.8 & Calculated\ocadhr
MACD versus signal & -0.65 and -0.55 & Calculated, 12, 26, 9\ocadhr
Average true range (14 days) & USD 0.98 & Calculated\ocadhr
20-day and 50-day average price & USD 38.75 and 40.41 & Calculated\ocadhr
200-day average price & USD 40.04 & Published, as of 23 October 2026\ocadhr
Volume, last 20 sessions versus 100 & 402,000 versus 325,000 & Calculated\ocadhr
\end{ocadtable}
\ocadtablenote{Source: Alpha Vantage daily series, 10 June to 30 October 2026; StockAnalysis.com. A benchmark index series was not retrieved, so market beta is the published figure.}
\ocadchart{relative.png}{BCH versus Banco Santander-Chile, rebased, 10 June to 30 October 2026}{Source: Alpha Vantage daily series. Units: index, 10 June 2026 = 100. BCH -5.4 percent; BSAC -2.4 percent.}
The technical structure is bearish in the short term: the price is below its 20, 50 and 200-day averages, the MACD is below its signal line and the RSI of 29.8 is at the oversold boundary, after a 15.1 percent decline from the August high. Volume has risen on the declines (20-day average 402,000 against 325,000). The very low published beta does not mean low risk; it reflects a local, peso-driven stock that moves with the Chilean market and with BSAC (correlation 0.88). Technical readings describe momentum, not value, and a 4.5 percent drop on 30 October has no cause that the retrieved sources identify.

\section{C. Financial Statements}
\subsection{C1. Income Statement}
\begin{ocadtable}{Income statement summary, CLP bn}{>{\raggedright\arraybackslash}p{50mm} R R R R R}
\ocadtablehead{\ocadH{Line}&\ocadH{FY2021}&\ocadH{FY2022}&\ocadH{FY2023}&\ocadH{FY2024}&\ocadH{FY2025}}
Interest income & 2,388.5 & 4,466.9 & 3,181.6 & 2,920.0 & 2,715.8\ocadhr
Interest expense & 814.4 & 2,200.8 & 1,634.7 & 1,138.3 & 969.6\ocadhr
Net interest income & 1,574.0 & 2,266.1 & 1,546.9 & 1,781.7 & 1,746.1\ocadhr
Non-interest income & 650.8 & 876.5 & 1,448.0 & 1,268.6 & 1,279.9\ocadhr
Operating revenue before loan losses & 2,224.9 & 3,142.6 & 2,994.9 & 3,050.3 & 3,026.0\ocadhr
Provision for loan losses & 19.2 & 412.1 & 361.3 & 391.8 & 381.9\ocadhr
Non-interest expense & 871.7 & 995.3 & 1,114.9 & 1,131.2 & 1,129.8\ocadhr
Pre-tax income & 1,332.3 & 1,735.0 & 1,517.5 & 1,525.8 & 1,513.2\ocadhr
Income tax & 276.0 & 289.2 & 273.9 & 318.4 & 320.9\ocadhr
Net income & 1,056.3 & 1,445.8 & 1,243.6 & 1,207.4 & 1,192.3\ocadhr
\end{ocadtable}
\ocadtablenote{Source: Banco de Chile financial statements (IFRS) as compiled by StockAnalysis.com, fiscal years to 31 December. Reported values.}
\ocadchart{mix.png}{Operating revenue composition, FY2021 to FY2025}{Source: StockAnalysis.com. Units: CLP bn.}
Interest income peaked at CLP 4,467 bn in 2022 with interest expense at CLP 2,201 bn, and net interest income recovered in 2024 as funding costs fell faster than asset yields. Non-interest income more than doubled between 2021 and 2025 (CLP 651 bn to CLP 1,280 bn), supplying 42.3 percent of revenue; the largest line, other non-interest income of CLP 1,147 bn, is not split between fees and indexation effects in the retrieved sources, and first-half 2026 net fees were CLP 340 bn. Operating expenses grew 29.6 percent between 2021 and 2025 against revenue growth of 36.0 percent, so operating leverage was positive over four years but neutral since 2023. Margins have contracted from 59.9 to 50.0 percent pre-tax mainly because provisions rose from CLP 19 bn to CLP 382 bn after the post-pandemic release in 2021.

The first half of 2026 shows the sensitivity of the model. Operating income was CLP 1,670.9 bn (up 8.4 percent), expenses CLP 576.6 bn (up 2.6 percent), credit loss expense CLP 279.3 bn (up 49.7 percent) and net income CLP 659.2 bn (up 4.0 percent). First-quarter net income was CLP 268.6 bn (calculated). In the second quarter the net interest margin was 5.76 percent against 4.71 percent a year earlier and return on average capital reached 29.1 percent. Full year guidance of 4.6 percent for the margin implies about 4.3 percent in the second half if both halves are weighted equally (calculated).

\subsection{C2. Balance Sheet}
\begin{ocadtable}{Balance sheet summary, CLP bn at 31 December}{>{\raggedright\arraybackslash}p{50mm} R R R R R}
\ocadtablehead{\ocadH{Line}&\ocadH{2021}&\ocadH{2022}&\ocadH{2023}&\ocadH{2024}&\ocadH{2025}}
Cash and equivalents & 2,655.0 & 3,152.9 & 2,289.7 & 2,035.1 & 1,657.9\ocadhr
Total investments & 11,961.7 & 13,466.5 & 13,079.1 & 7,699.9 & 9,620.8\ocadhr
Net loans & 33,951.2 & 36,276.7 & 37,250.0 & 38,755.9 & 38,754.1\ocadhr
Total assets & 51,758.3 & 55,108.9 & 55,792.6 & 52,095.4 & 54,100.9\ocadhr
Total deposits & 27,053.6 & 27,951.1 & 28,687.2 & 28,432.0 & 28,470.2\ocadhr
Total debt (borrowings and bonds) & 17,580.4 & 19,306.6 & 18,720.3 & 15,011.5 & 16,211.8\ocadhr
Total liabilities & 47,464.8 & 49,677.4 & 50,555.3 & 46,472.4 & 48,301.4\ocadhr
Shareholders' equity & 4,293.5 & 5,431.4 & 5,237.3 & 5,623.0 & 5,799.5\ocadhr
Allowance for loan losses & 718.8 & 823.3 & 769.7 & 787.0 & 836.2\ocadhr
Non-interest bearing deposits, \% of deposits & 56.7 & 40.0 & 38.4 & 41.4 & 41.4\ocadhr
\end{ocadtable}
\ocadtablenote{Source: StockAnalysis.com balance sheets of Banco de Chile (reported); percentages calculated. Net debt, debt to EBITDA and current ratio are not meaningful for a bank.}
\ocadchart{bs.png}{Net loans, deposits and borrowings, 2021 to 2025}{Source: StockAnalysis.com. Units: CLP trillion.}
Net loans grew 14.1 percent between 2021 and 2025, deposits 5.2 percent, and wholesale debt filled the gap, with net loans at 136 percent of deposits at end 2025. This funding structure is common in Chile, where bonds and institutional funding supplement deposits, but it makes funding costs sensitive to market rates. The allowance covered 2.09 percent of gross loans at end 2025 (1.97 percent in 2024). Equity rose 35.1 percent over four years, lifting equity to assets from 8.3 to 10.7 percent. Tangible book value was CLP 5,625 bn and goodwill is nil. Regulatory capital, liquidity coverage and non-performing loan ratios are Not Retrieved for 2025; the past-due ratio was 1.65 percent in June 2026 against 1.47 percent a year earlier according to the 2Q26 filing summary.

\subsection{C3. Cash Flow Statement}
\begin{ocadtable}{Cash flow summary, CLP bn}{>{\raggedright\arraybackslash}p{50mm} R R R R R}
\ocadtablehead{\ocadH{Line}&\ocadH{FY2021}&\ocadH{FY2022}&\ocadH{FY2023}&\ocadH{FY2024}&\ocadH{FY2025}}
Net income & 1,056.3 & 1,445.8 & 1,243.6 & 1,207.4 & 1,192.3\ocadhr
Operating cash flow & -598.2 & 51.0 & -186.7 & 396.6 & 1,248.4\ocadhr
Capital expenditure & -35.6 & -21.2 & -26.7 & -17.2 & -18.7\ocadhr
Operating cash flow less capex & -633.8 & 29.8 & -213.4 & 379.4 & 1,229.7\ocadhr
Common dividends paid & -220.3 & -539.8 & -866.9 & -815.9 & -995.4\ocadhr
Net debt issued (repaid) & 1,596.7 & 378.0 & -674.0 & -516.4 & 668.0\ocadhr
Change in deposits & 3,603.8 & 77.6 & 1,228.1 & -225.3 & 54.0\ocadhr
\end{ocadtable}
\ocadtablenote{Source: StockAnalysis.com cash flow statements (reported). Operating cash flow less capex equals the provider's free cash flow; for a bank it is not an indicator of distributable cash.}
For a bank the cash flow statement is dominated by changes in trading securities and loans rather than by earnings: in 2025 trading assets absorbed CLP 1,699.6 bn while other operating assets released CLP 712.4 bn. Net income therefore does not equal cash generation, and operating cash flow swung from negative CLP 598 bn in 2021 to positive CLP 1,248 bn in 2025. Capital expenditure is small at CLP 19 bn (0.6 percent of revenue), and dividends paid rose from CLP 220 bn to CLP 995 bn, which is 83 percent of 2025 net income on a cash basis; this includes timing effects because the published payout ratio is 67.5 percent. Earnings are supported by the balance sheet rather than by free cash flow, and the retained portion rose equity.

\subsection{C4. Forecasting and Projections}
Revenue is driven by volumes and the margin. Loan growth is guided near 6 percent nominal for 2026 (lowered from 7 percent) against 3.8 percent for the system, implying share gains; the net interest margin is guided to 4.6 percent, the efficiency ratio to 37.0 percent, credit cost to 1.2 to 1.3 percent of average loans, and return on average capital to 21 to 22 percent. Working capital and capex are immaterial for a bank; growth capex is mainly technology, with CLP 59 bn of intangible purchases in 2025.
\begin{ocadtable}{Net income scenarios, CLP bn}{>{\raggedright\arraybackslash}p{34mm} R R R L}
\ocadtablehead{\ocadH{Scenario}&\ocadH{FY2026E}&\ocadH{FY2027E}&\ocadH{FY2028E}&\ocadH{Key assumptions}}
Bear & 1,150 & 1,100 & 1,090 & Margin falls faster, credit cost 1.6 percent of loans\ocadhr
Base & 1,260 & 1,300 & 1,345 & Guidance met, growth of 3 to 3.5 percent\ocadhr
Bull & 1,320 & 1,380 & 1,450 & Margin holds near 4.9 percent, credit cost 1.2 percent\ocadhr
\end{ocadtable}
\ocadtablenote{Source: analyst estimates built from the 2026 guidance in the Banco de Chile 2Q26 filing. FY2025 reported net income was CLP 1,192 bn. Base 2026 equals 1H26 reported CLP 659 bn plus about CLP 601 bn in the second half.}
\ocadchart{scen.png}{Net income scenarios, FY2026E to FY2028E}{Source: analyst estimates. Units: CLP bn.}
The base case equals about CLP 2,495 per ADS (CLP 1,260 bn over 505.1 million ADSs), a forward P/E of about 14 times at the current price using an implied exchange rate near CLP 970 per USD. This rate is inferred from the published price to book and is not a quoted market rate. Sensitivity of valuation to margin and credit cost is shown in the next section through the return on equity.

\subsection{C5. Discounted Cash Flow (DCF) Model}
A free cash flow DCF is not appropriate for a bank, because debt is operating raw material and operating cash flow is distorted by trading assets. The standard approach is an excess return (dividend discount) model, in which value equals book value multiplied by (ROE minus g) divided by (cost of equity minus g). Assumptions: book value per ADS of USD 11.85 (implied by the published price to book); sustainable ROE of 19 percent (below the 2026 guidance of 21 to 22 percent, reflecting normalization of inflation); a payout consistent with growth of 2 to 4 percent; tax of about 21 percent and no separate capex or working capital forecast. The cost of equity range is 9 to 13 percent; net debt is not applicable and the share count is unchanged at 101,017 million (505.1 million ADSs).
\begin{ocadtable}{Sensitivity: cost of equity (rows) against terminal growth (columns), value per ADS, USD}{L R R R}
\ocadtablehead{\ocadH{Cost of equity}&\ocadH{g = 2\%}&\ocadH{g = 3\%}&\ocadH{g = 4\%}}
9\% & 28.8 & 31.6 & 35.6\ocadhr
10\% & 25.2 & 27.1 & 29.6\ocadhr
11\% & 22.4 & 23.7 & 25.4\ocadhr
12\% & 20.1 & 21.1 & 22.2\ocadhr
13\% & 18.3 & 19.0 & 19.8\ocadhr
\end{ocadtable}
\ocadtablenote{Source: analyst calculation. Sustainable ROE of 19 percent and book value of USD 11.85 per ADS. With a cost of equity of 10 percent and growth of 3 percent, each 2 points of sustainable ROE adds about USD 3.4 per ADS (15 percent ROE gives USD 20.3 and 23 percent gives USD 33.9).}
The calculation gives a range of USD 18 to 36 per ADS, with most of the range between USD 21 and USD 30 at 10 to 12 percent cost of equity. The current price of USD 36.55 sits at or above the top of the range and is reached only at a 9 percent cost of equity with growth of 4 percent or a sustainable ROE above 23 percent. This is a valuation range under stated assumptions, not a price target or a forecast.

\subsection{C6. Leveraged Buyout (LBO) Model}
Not applicable. A bank is financed by regulated deposit funding and capital requirements, the controlling holders hold 51.15 percent, and a leveraged buyout structure is not meaningful.

\section{D. Investment Thesis}
\subsection{Bull Case}
Inflation stays near 4 percent, the margin holds close to 4.9 percent rather than falling to 4.3 percent in the second half, credit cost stays at 1.2 percent of loans, and loan growth of 6 percent with a stable 37 percent efficiency ratio lifts net income towards CLP 1,380 bn in 2027. A return on equity above 21 percent would justify a price to book near 3 times at a 10 percent cost of equity.
\subsection{Bear Case}
Margin falls to the guided level or below, provisions rise to 1.6 percent of loans, and growth stalls at the system's 3.8 percent nominal rate. Net income would fall towards CLP 1,100 bn in 2027, return on equity would approach 18 percent, and on the excess return model a value near USD 20 to 25 per ADS compares with USD 36.55.
\subsection{Base Case}
Guidance is met: net income of about CLP 1,260 bn in 2026, rising 3 to 3.5 percent a year, return on equity of 20 to 21 percent, a payout near two thirds and a dividend yield near 4.5 percent. At this earnings level the stock trades at about 14 times 2026 estimated earnings.
\subsection{Catalysts}
\begin{ocadtable}{Potential catalysts}{>{\raggedright\arraybackslash}p{36mm} >{\raggedright\arraybackslash}p{26mm} L}
\ocadtablehead{\ocadH{Event}&\ocadH{Timing}&\ocadH{Why it matters}}
3Q26 results release & Per company calendar; not retrieved & Tests the margin guidance and credit cost trend.\ocadhr
Dividend decision & Around March 2027, based on the 2026 pattern & Ex-dividend date was 30 March 2026; payout sets the yield.\ocadhr
Guidance update & Quarterly & Loan growth already cut once from 7 to 6 percent.\ocadhr
Central Bank policy rate & Scheduled meetings & Drives funding costs and UF indexation.\ocadhr
Chile inflation and GDP data & Monthly and quarterly & Sets nominal growth and arrears.\ocadhr
\end{ocadtable}
\subsection{Things to Monitor}
\begin{ocadtable}{Monitoring list}{>{\raggedright\arraybackslash}p{34mm} L L L}
\ocadtablehead{\ocadH{Metric or event}&\ocadH{Why it matters}&\ocadH{Positive}&\ocadH{Negative}}
Net interest margin & Largest earnings driver & Stays above 4.6 percent & Falls below 4.3 percent\ocadhr
Credit loss expense & Earnings volatility & At or below 1.3 percent of loans & Above 1.6 percent\ocadhr
Past-due ratio & Early warning on credit & Below 1.65 percent & Above 2.0 percent\ocadhr
Efficiency ratio & Operating leverage & Near or below 37 percent & Above 39 percent\ocadhr
Return on average capital & Valuation support & 21 to 22 percent or higher & Below 18 percent\ocadhr
Non-interest bearing deposits & Funding cost & Above 41 percent of deposits & Below 38 percent\ocadhr
\end{ocadtable}

\section{E. Final Investment Framework}
\begin{ocadframework}
\fwrow{Business quality}{High. A universal franchise with 2.9 million clients and a low-cost deposit base.}
\fwrow{Financial quality}{High but cyclical. Return on equity of 20.9 percent and efficiency of 37.3 percent in 2025, with net income down in each year since 2022.}
\fwrow{Competitive position}{Strong. Second-largest private bank by loans, approximately 15 to 16 percent of system loans (indicative, comparing end 2025 bank loans to May 2026 system loans).}
\fwrow{Management}{Competent and shareholder-friendly. Stable share count, steady payout, no goodwill; controlling holders own 51.15 percent.}
\fwrow{Industry}{Mature, cyclical and regulated; loan growth near 3.8 percent nominal.}
\fwrow{Valuation}{Full. 3.1 times book and 15.2 times earnings, with the price above the excess return range of USD 21 to 30.}
\fwrow{Technical and market structure}{Weak near term. Below the 20, 50 and 200-day averages, RSI 29.8, down 15.1 percent from the August high.}
\fwrow{Principal risk}{Normalization of the net interest margin as inflation and policy rates decline.}
\fwrow{Principal catalyst}{Third-quarter results confirming or revising the 2026 margin and credit cost guidance.}
\fwrow{Key variable}{Sustainable return on equity: each 2 points changes the excess return value by about USD 3.4 per ADS.}
\fwrow{Thesis invalidation}{The investment thesis would become materially weaker if the net interest margin falls below 4.3 percent, credit loss expense exceeds 1.6 percent of average loans, the efficiency ratio rises above 39 percent, or return on average capital falls below 18 percent for two consecutive quarters.}
\end{ocadframework}

\section{F. Data Table}
\ocadlongtablesetup
\begin{longtable}{@{}p{18mm}p{40mm}p{28mm}p{28mm}p{28mm}@{}}
\rowcolor{ocadNavy}\ocadH{Category}&\ocadH{Metric}&\ocadH{Current}&\ocadH{Prior year}&\ocadH{3Y / 5Y trend}\ocadhr
\endfirsthead
\rowcolor{ocadNavy}\ocadH{Category}&\ocadH{Metric}&\ocadH{Current}&\ocadH{Prior year}&\ocadH{3Y / 5Y trend}\ocadhr
\endhead
Market & Market cap, USD bn & 18.5 & 19.2 (Dec 2025) & 7.9 (Dec 2021)\ocadhr
Market & Enterprise value & N/A & N/A & N/A\ocadhr
Market & Share price, USD per ADS & 36.55 & 37.99 (Dec 2025) & 15.70 (Dec 2021)\ocadhr
Growth & Revenue growth FY2025, \% & -0.8 & 1.8 & CAGR 8.0 (2021 to 2025)\ocadhr
Profitability & Gross margin & N/A & N/A & N/A\ocadhr
Profitability & Pre-tax margin, \% & 50.0 & 50.0 & 59.9 (FY2021)\ocadhr
Profitability & Net margin, \% & 39.4 & 39.6 & 47.5 (FY2021)\ocadhr
Profitability & ROE, \% & 20.9 & 22.2 & 23.3 (FY2023)\ocadhr
Profitability & ROIC & N/A & N/A & N/A\ocadhr
Cash flow & Operating cash flow, CLP bn & 1,248.4 & 396.6 & -186.7 (FY2023)\ocadhr
Cash flow & Operating cash flow less capex, CLP bn & 1,229.7 & 379.4 & -213.4 (FY2023)\ocadhr
Leverage & Net debt / EBITDA & N/A & N/A & N/A\ocadhr
Leverage & Equity to assets, \% & 10.7 & 10.8 & 9.4 (FY2023)\ocadhr
Valuation & P/E, x & 15.2 & Not Retrieved & Not Retrieved\ocadhr
Valuation & P/B, x & 3.08 & Not Retrieved & Not Retrieved\ocadhr
Valuation & EV/EBITDA & N/A & N/A & N/A\ocadhr
Valuation & P/FCF & N/A & N/A & N/A\ocadhr
Market & Beta (five years) & 0.06 & Not Retrieved & Not Retrieved\ocadhr
Market & Volatility, annualized \% & 25.5 & Not Retrieved & Not Retrieved\ocadhr
Market & Maximum drawdown, \% & -15.1 (100 sessions) & Not Retrieved & Not Retrieved\ocadhr
\end{longtable}
\ocadtablenote{Source: StockAnalysis.com; Alpha Vantage; calculations by the analyst. N/A means not meaningful for a bank; Not Retrieved means unavailable at the report date.}

\section{G. Sources and Methodology}
\subsection{Data cross-check}
Financial statements were taken from the compiled annual statements and cross-checked with the company page of the controlling shareholder: net income of CLP 1,192,262 million, total assets of CLP 54,100,903 million and equity of CLP 5,799,534 million match exactly. The reported ROAE of 21.2 percent compares with 20.9 percent calculated on average equity, and the reported efficiency ratio of 37.4 percent with 37.3 percent calculated, which is a definition difference. Reported total loans of CLP 39,191,970 million differ from the compiled gross loans of CLP 39,917,824 million for the same reason, and the reported figure is used for market share. The closing price of USD 36.55 from the daily series agrees with the market capitalization of USD 18.46 bn divided by 505.085 million ADSs. The payout ratio is published as 63.9 percent on one page and 67.5 percent on another; the higher figure from the statistics page is used and the discrepancy noted. A book value per share shown as 0.06 on the statistics page is a data error and was ignored. Third-party multiples dated 23 October were rescaled by the change in price to 30 October. All figures are IFRS, in Chilean pesos unless marked USD, and fiscal years end on 31 December.
\subsection{Source list}
\ocadsourcegroup{Primary Sources}
\ocadsource{Banco de Chile}{Form 6-K, Exhibit 99.1, 2Q26 Financial Management Review (copy hosted by StockTitan)}{31 July 2026}{https://www.stocktitan.net/sec-filings/BCH/6-k-bank-of-chile-current-report-foreign-issuer-b0b7e2580c31.html}{Second quarter and first half 2026 results, ownership of 51.15 percent, 2026 guidance, industry loans.}
\ocadsource{Banco de Chile}{Announcement of Annual Report on Form 20-F 2025}{30 April 2026}{https://www.accessnewswire.com/newsroom/en/banking-and-financial-services/banco-de-chile-announces-annual-report-form-20-f-2025-1162983}{Filing of the FY2025 annual report.}
\ocadsource{Quiñenco S.A.}{Banco de Chile company page}{31 December 2025 data}{https://www.quinenco.cl/en/banco-de-chile-company/}{Ownership through LQIF, Citigroup partnership, clients, employees, FY2025 reported results and ratios.}
\ocadsource{Banco Santander-Chile}{Form 20-F for fiscal year 2025}{2026}{https://www.sec.gov/Archives/edgar/data/1027552/000102755226000038/bsac-20251231.htm}{Chile GDP, inflation and policy rate history; peer loan mix.}
\ocadsourcegroup{Secondary Sources}
\ocadsource{StockAnalysis.com}{Banco de Chile income statement}{FY2025 and earlier}{https://stockanalysis.com/stocks/bch/financials/income-statement/}{Annual income statement FY2021 to FY2025.}
\ocadsource{StockAnalysis.com}{Banco de Chile balance sheet}{FY2025 and earlier}{https://stockanalysis.com/stocks/bch/financials/balance-sheet/}{Annual balance sheet FY2021 to FY2025.}
\ocadsource{StockAnalysis.com}{Banco de Chile cash flow statement}{FY2025 and earlier}{https://stockanalysis.com/stocks/bch/financials/cash-flow-statement/}{Annual cash flow FY2021 to FY2025.}
\ocadsource{StockAnalysis.com}{Banco de Chile company profile}{October 2026}{https://stockanalysis.com/stocks/bch/company/}{Executive team and employee count.}
\ocadsource{StockAnalysis.com}{Banco de Chile statistics}{23 October 2026}{https://stockanalysis.com/stocks/bch/statistics/}{Multiples, ROE, beta, moving averages, payout.}
\ocadsource{StockAnalysis.com}{Banco Santander-Chile statistics}{24 September 2026}{https://stockanalysis.com/stocks/bsac/statistics/}{Peer multiples.}
\ocadsource{StockAnalysis.com}{Bancolombia statistics}{26 October 2026}{https://stockanalysis.com/stocks/cib/statistics/}{Regional peer multiples.}
\ocadsource{StockAnalysis.com}{Credicorp statistics}{October 2026}{https://stockanalysis.com/stocks/bap/statistics/}{Regional peer multiples.}
\ocadsourcegroup{Market Data}
\ocadsource{Alpha Vantage}{Daily time series, BCH and BSAC}{30 October 2026}{https://www.alphavantage.co/documentation/}{Daily open, high, low, close and volume, 10 June to 30 October 2026.}
\ocadsource{StockAnalysis.com}{Banco de Chile market capitalization history}{30 October 2026}{https://stockanalysis.com/stocks/bch/market-cap/}{Year-end market capitalization 2002 to 2025 and current.}
\ocadsource{StockAnalysis.com}{Banco de Chile dividend history}{April 2026}{https://stockanalysis.com/stocks/bch/dividend/}{ADS dividends paid in 2024, 2025 and 2026.}
\ocadsourcegroup{Calculations}
Ratios in tables marked calculated were derived by the analyst: return on average equity and assets use the mean of opening and closing balances; efficiency ratio is non-interest expense divided by operating revenue before loan losses; provisions to gross loans uses year-end gross loans; returns are close to close, excluding dividends; volatility is the standard deviation of 99 daily log returns times the square root of 252; RSI, MACD and ATR use standard 14, 12-26-9 and 14-day settings. Year-end share prices are year-end market capitalization divided by 505.085 million ADSs (101,017 million shares over 200). The justified price to book model uses the formula (ROE minus g) over (cost of equity minus g); the cost of equity range and sustainable ROE are analyst assumptions. Implied exchange rate near CLP 970 per USD is inferred from the published price to book. Market share of about 15 to 16 percent is CLP 39,191,970 million over CLP 253,069,875 million and is indicative because the dates differ.

\begin{ocaddisclaimer}
This report has been prepared for personal research purposes only and does not constitute investment advice, financial advice, or a recommendation to buy, sell, or hold any security. The information contained herein has been compiled with the assistance of artificial intelligence and is based on publicly available sources believed to be reliable, but its accuracy and completeness are not guaranteed. Past performance is not indicative of future results. Investors should conduct their own due diligence and consult with a qualified financial professional before making any investment decisions.
\end{ocaddisclaimer}
\end{document}
